\documentclass[10pt,a4paper]{amsart}
\usepackage[margin=19mm]{geometry}
\usepackage{amsmath,amssymb,mathtools}
\usepackage[scr=rsfs,cal=euler]{mathalfa}
\usepackage{xfrac}
\usepackage[unicode,hidelinks,bookmarks=false]{hyperref}
\newcommand{\ie}{i.e.\ }
\newcommand{\cf}{cf.\ }
\newcommand{\resp}{respectively\ }
\newcommand{\Subsection}{\subsection}
\providecommand{\nobreakdash}{\nobreak\hbox{-}}
\setlength{\emergencystretch}{2em}
\multlinegap0pt
\usepackage{amssymb}
\usepackage{tikz}
\newtheorem{lemma}{Lemma}
\title{Isoperimetric Inequality for degenerate elliptic operators of Grushin type}
\author{Dangyang He}
\date{Source: arXiv:2605.08818v1 (CC BY 4.0)}
\begin{document}
\maketitle

\noindent\emph{Source and licence.} Isoperimetric Inequality for degenerate elliptic operators of Grushin type. Version arXiv:2605.08818v1 (CC BY 4.0), distributed by arXiv under the Creative Commons Attribution 4.0 International licence, http://creativecommons.org/licenses/by/4.0/, as recorded in the arXiv licence metadata for this exact version and shown on the abstract page https://arxiv.org/abs/2605.08818v1. Full licence notice: \texttt{licenses/arxiv-2605.08818-CC-BY-4.0.txt}.\par\smallskip

\noindent\textit{Editorial scope.} Counted unit: the article's lemma le\_Ric\_Grushin (source0.tex lines 402--408). The article's complete original proof of it is delivered with it (source0.tex lines 410--517, one \texttt{proof} environment of 108 lines). No mathematics is written, repaired or re-derived: the statement and the proof are the article's own lines, in the article's own notation, and no retained line is substituted or re-flowed.  Retained uncounted as the source-local context the proof needs: lines 389--391; lines 393--395.  Nothing was omitted from any retained span, so the package carries no omission record.  No retained line contains a citation, so the wrapper carries no bibliography and no bibliography entry.  No retained line contains a figure environment, an image-inclusion command or drawing code, so no image is delivered and the package declares no asset dependency.  Editorial formatting only, in addition: an AMS \texttt{amsart} preamble that restates the article's own declarations and macros used by the retained lines, namely the environments \texttt{lemma} on separate counters, together with the standard packages those lines need.  The retained environments print in this document as Lemma 1 in order of appearance, whereas the article prints the same items with the numbering of its own full document.  The complete native source of the article as distributed in the arXiv e-print for this exact version is bundled unchanged as \texttt{sources/200177/source0.tex}.
\begin{align}\label{carre1}
    \Gamma(f,g) = \frac{1}{2} \left( \mathcal{L}(fg) - f \mathcal{L}(g) - g\mathcal{L}(f)\right).
\end{align}

\begin{align}\label{carre2}
    \Gamma_2(f,g) = \frac{1}{2} \left( \mathcal{L}\Gamma(f,g) - \Gamma(f,\mathcal{L}(g)) - \Gamma(g,\mathcal{L}(f)) \right),
\end{align}

\begin{lemma}\label{le_Ric_Grushin}
Let $n,m\ge 1$, $\alpha \in (0,1)$ and $\beta \ge 0$. Let $\Gamma$ (resp. $\Gamma_2$) be the \textit{carre du champ} (resp. second \textit{carre du champ}) of $L$. The following curvature dimension inequality holds
\begin{align*}
\Gamma_2(f) \ge - \frac{C(n,m,\alpha,\beta)}{|x|^{2(1-\alpha)}} \Gamma(f) + C'(n,m,\alpha,\beta) \left(L(f) \right)^2
\end{align*}
for all $f\in C^\infty( \mathbb{R}^{n}\setminus\{0\} \times \mathbb{R}^m)$.
\end{lemma}

\begin{proof}[Proof of Lemma~\ref{le_Ric_Grushin}]
Since $L$ is non-negative, we set $\mathcal{L} = -L$ and calculate \eqref{carre1} and \eqref{carre2} directly. Put $r=|x|$, $a = r^{2\alpha}$ and $b=r^{2\beta}$. Then
\begin{equation*}
\mathcal L = -a\Delta_x - b\Delta_y+\nabla a\cdot \nabla_x.
\end{equation*}

For simplicity, we use indices $i,j\in\{1,\dots,n\}$, $\mu,\nu\in\{1,\dots,m\}$, one gets directly from the definition \eqref{carre1} that
\begin{equation}\label{eq_gamma1}
    \Gamma(f)=a f_i^2+b f_\mu^2.
\end{equation}

Next, by \eqref{carre2}, a direct calculation gives
\begin{align}\label{eq_gamma2}
\Gamma_2(f)
={}&a^2 f_{ij}^2+2ab f_{i\mu}^2+b^2 f_{\mu\nu}^2\\ \nonumber
&+2a a_i f_j f_{ij} + a a_j f_j \Delta_x f\\ \nonumber
&+2a b_i f_\mu f_{i\mu}+ a b_i f_i \Delta_y f\\ \nonumber
&+\Bigl(-\frac a2 \Delta_x a+\frac12|\nabla a|^2\Bigr)f_i^2
-a a_{ij}f_i f_j\\ \nonumber
&+\Bigl(-\frac a2 \Delta_x b+\frac12\nabla a\cdot \nabla b\Bigr)f_\mu^2.
\end{align}
Note that
\begin{align*}
    a_i=2\alpha r^{2\alpha-2}x_i,\quad
|\nabla a|^2 = 4\alpha^2 r^{4\alpha-2},
\quad
\Delta_x a = -2\alpha(n+2\alpha-2)r^{2\alpha-2},
\end{align*}
\begin{equation*}
    a_{ij} = 2\alpha r^{2\alpha-2}\delta_{ij}
 + 4\alpha(\alpha-1)r^{2\alpha-4}x_i x_j,
\end{equation*}
and 
\begin{align*}
    b_i = 2\beta r^{2\beta-2}x_i,\quad
|\nabla b|^2 = 4\beta^2 r^{4\beta-2},
\quad
\Delta_x b = -2\beta(n+2\beta-2)r^{2\beta-2}.
\end{align*}


We estimate the lower-order terms in \eqref{eq_gamma2} by Young's inequality.

First, using $|\Delta_x f|\le \sqrt n|D_x^2 f|$, where $|D_x^2f|^2 = \sum_{1\le i,j\le n}|\partial_{x_i,x_j}f|^2$, and Cauchy-Schwartz inequality and then Young's inequality, one obtains
\begin{equation}\label{1}
    2a a_i f_j f_{ij} + a a_j f_j \Delta_x f \ge
-\frac14 a^2|D_x^2 f|^2-C|\nabla a|^2|\nabla_x f|^2.
\end{equation}
Also, since $\alpha<1$,
\begin{equation}\label{2}
-a a_{ij}f_i f_j = -2\alpha r^{4\alpha-2}|\nabla_x f|^2 + 4\alpha(1-\alpha)r^{4\alpha-4}(x\cdot \nabla_x f)^2
\ge -2\alpha r^{4\alpha-2}|\nabla_x f|^2.
\end{equation}
Hence
\begin{equation}\label{3}
    \Bigl(-\frac a2 \Delta_x a + \frac12| \nabla a|^2\Bigr)f_i^2-a a_{ij}f_i f_j \ge
-C r^{4\alpha-2}|\nabla_x f|^2.
\end{equation}

Next, it is clear that
\begin{equation}\label{4}
    2a b_i f_\mu f_{i\mu}
\ge
-a b f_{i\mu}^2-\frac{a|\nabla b|^2}{b} |\nabla_y f|^2,
\end{equation}
and, using $|\Delta_y f|\le \sqrt m|D_y^2 f|$, where $|D_y^2f|^2 = \sum_{1\le \mu,\nu\le m}|\partial_{y_\mu,y_\nu}f|^2$, 
\begin{equation}\label{5}
    a b_i f_i \Delta_y f \ge -\frac14 b^2|D_y^2 f|^2 - C \frac{a^2|\nabla b|^2}{b^2}|\nabla_x f|^2.
\end{equation}

Finally,
\begin{equation}\label{6}
    \Bigl( -\frac a2 \Delta_x b + \frac12 \nabla a \cdot \nabla b \Bigr) f_\mu^2 \ge
 -  C r^{2\alpha+2\beta-2}|\nabla_y f|^2.
\end{equation}

Insert \eqref{1}--\eqref{6} into \eqref{eq_gamma2}. 
% \begin{equation*}
%     |\nabla a|^2\sim r^{4\alpha-2},\quad
% \frac{a|\nabla b|^2}{b}\sim r^{2\alpha+2\beta-2},\quad
% \frac{a^2|\nabla b|^2}{b^2}\sim r^{4\alpha-2},
% \end{equation*}
We obtain
\begin{align}\label{7}
{}&\Gamma_2(f)\\ \nonumber
&\ge \frac34 a^2|D_x^2 f|^2 + ab|D_{xy}^2 f|^2 + \frac34 b^2|D_y^2 f|^2 - C \Bigl(r^{4\alpha-2}|\nabla_x f|^2 + r^{2\alpha+2\beta-2}|\nabla_y f|^2\Bigr)\\ \nonumber
&= \frac34 a^2|D_x^2 f|^2 + ab|D_{xy}^2 f|^2 + \frac34 b^2|D_y^2 f|^2 -\frac{C}{|x|^{2-2\alpha}}\Gamma(f)
\end{align}

On the other hand, by trace inequality, with $C'=12(n+m)$, 
\begin{align}\label{8}
C'^{-1}(\mathcal L f)^2 &\le
C'^{-1} 3a^2(\Delta_x f)^2+ C'^{-1} 3b^2(\Delta_y f)^2+ C'^{-1}3|\nabla a|^2|\nabla_x f|^2\\ \nonumber
&\le C'^{-1} 3n a^2|D_x^2 f|^2 + C'^{-1} 3m b^2|D_y^2 f|^2 + C'^{-1} 12\alpha^2 r^{4\alpha-2}|\nabla_x f|^2\\ \nonumber
&\le \frac14 a^2|D_x^2 f|^2+\frac14 b^2|D_y^2 f|^2+\frac1{|x|^{2-2\alpha}}\Gamma(f).
\end{align}

Combining \eqref{7} and \eqref{8}, we conclude that
\begin{equation*}
    \Gamma_2(f)\ge -\frac{C}{|x|^{2-2\alpha}}\Gamma(f)+\frac1{C'}(\mathcal L f)^2
\end{equation*}
for some $C=C(n,m,\alpha,\beta)>0$, as desired.





\end{proof}

\end{document}
